\documentclass[tikz,border=14pt]{standalone}
\usepackage[utf8]{inputenc}
\usepackage[T1]{fontenc}
\usepackage{vntex}
\usetikzlibrary{calc}

% ==== Bảng màu ====
\definecolor{hw1}{HTML}{3B5BA5}
\definecolor{hw2}{HTML}{6A55B2}
\definecolor{hw3}{HTML}{B4418E}
\definecolor{hw4}{HTML}{D6336C}
\definecolor{cT1}{HTML}{6FA8DC}
\definecolor{cT2}{HTML}{5B8DD6}
\definecolor{cT3}{HTML}{9B6FD6}
\definecolor{cT4}{HTML}{C86DD7}
\definecolor{cT5}{HTML}{E94E89}
\definecolor{cT6}{HTML}{4FB3A9}
\definecolor{cMS}{HTML}{C0392B}

% Thứ trong tuần (0 = Thứ Hai)
\newcommand{\dayname}[1]{\ifcase#1 T2\or T3\or T4\or T5\or T6\or T7\or CN\fi}

\begin{document}
\begin{tikzpicture}[x=0.62cm, y=0.62cm, font=\sffamily]

% Trục thời gian: cột d (0..27) = ngày 07/09/2026 + d (07/09 là Thứ Hai)
\def\lw{18}     % độ rộng cột nhãn (đơn vị x)
\def\nd{28}     % 4 tuần
\def\nr{34}     % số hàng nội dung

% ==== Tiêu đề ====
\node[anchor=west,font=\sffamily\bfseries\large] at (-\lw,3.0)
  {Kế hoạch thực hiện Thực tập 1 -- Tổng hợp \& tái lập bài báo FashionMV (ProCIR)};
\node[anchor=west,font=\sffamily\small] at (-\lw,2.15)
  {Học viên: Võ Thị Bích Phượng -- 2470570 \qquad GVHD: PGS.TS. Võ Thị Ngọc Châu};

% ==== Tô nền cuối tuần ====
\foreach \d in {0,...,27} {
  \pgfmathtruncatemacro{\k}{mod(\d,7)}
  \ifnum\k>4 \fill[gray!9] (\d,0) rectangle (\d+1,-\nr); \fi
}

% ==== Header tuần ====
\newcommand{\weekhead}[3]{%
  \fill[#3] (#1*7,0.75) rectangle (#1*7+7,1.55);
  \node[white,font=\sffamily\bfseries\scriptsize] at (#1*7+3.5,1.15) {#2};
}
\weekhead{0}{Tuần 1: 07/09 -- 13/09}{hw1}
\weekhead{1}{Tuần 2: 14/09 -- 20/09}{hw2}
\weekhead{2}{Tuần 3: 21/09 -- 27/09}{hw3}
\weekhead{3}{Tuần 4: 28/09 -- 04/10}{hw4}

% ==== Header ngày ====
\foreach \d in {0,...,27} {
  \pgfmathtruncatemacro{\k}{mod(\d,7)}
  \node[font=\sffamily\tiny] at (\d+0.5,0.375) {\dayname{\k}};
}
\draw[black] (0,0) rectangle (\nd,1.55);
\draw[black] (0,0.75) -- (\nd,0.75);

% ==== Lưới ====
\foreach \d in {0,...,28} \draw[gray!25] (\d,0) -- (\d,-\nr);
\foreach \r in {0,...,34} \draw[gray!25] (-\lw,-\r) -- (\nd,-\r);
\foreach \w in {0,...,4} \draw[gray!70] (\w*7,0.75) -- (\w*7,-\nr);
\draw[black] (0,0) rectangle (\nd,-\nr);

% ==== Lệnh vẽ ====
\newcommand{\taskrow}[2]{%
  \fill[black] (-\lw,-#1) rectangle (\nd,-#1-1);
  \node[white,anchor=west,font=\sffamily\bfseries\small] at (-\lw+0.3,-#1-0.5) {#2};
}
\newcommand{\act}[2]{%
  \node[anchor=west,font=\sffamily\footnotesize] at (-\lw+0.9,-#1-0.5) {#2};
}
% \gbar{hàng}{ngày bắt đầu}{ngày kết thúc (không gồm)}{màu}
\newcommand{\gbar}[4]{%
  \fill[#4,rounded corners=1.5pt] (#2+0.06,-#1-0.18) rectangle (#3-0.06,-#1-0.82);
}
% mốc (hình thoi) tại giữa ngày d
\newcommand{\milestone}[2]{%
  \fill[cMS] (#2+0.5,-#1-0.1) -- (#2+0.9,-#1-0.5) -- (#2+0.5,-#1-0.9) -- (#2+0.1,-#1-0.5) -- cycle;
}

% ========================================================
\taskrow{0}{TASK 1 --- Khảo sát tài liệu}
\act{1}{Tìm hiểu bài toán Composed Image Retrieval (CIR)}
\gbar{1}{0}{2}{cT1}
\act{2}{Khảo sát các hướng tiếp cận CIR (fusion, CLIP, zero-shot, MLLM)}
\gbar{2}{1}{4}{cT1}
\act{3}{Khảo sát các bộ dữ liệu CIR (FashionIQ, CIRR, FACap)}
\gbar{3}{3}{5}{cT1}
\act{4}{Tổng hợp tổng quan tài liệu (literature review)}
\gbar{4}{4}{7}{cT1}

% ========================================================
\taskrow{5}{TASK 2 --- Nghiên cứu bài báo FashionMV}
\act{6}{Bài toán Multi-View CIR và hiện tượng View Incompleteness}
\gbar{6}{5}{7}{cT2}
\act{7}{Pipeline xây dựng bộ dữ liệu FashionMV}
\gbar{7}{6}{8}{cT2}
\act{8}{Kiến trúc ProCIR (hội thoại hai lượt, Align, CoT, SFT)}
\gbar{8}{7}{10}{cT2}
\act{9}{Kết quả thực nghiệm và phân tích ablation}
\gbar{9}{9}{11}{cT2}

% ========================================================
\taskrow{10}{TASK 3 --- Đề cương nghiên cứu}
\act{11}{Viết đề cương nghiên cứu (bản nháp)}
\gbar{11}{10}{13}{cT3}
\act{12}{Chỉnh sửa đề cương theo góp ý GVHD}
\gbar{12}{13}{16}{cT3}
\act{13}{Chốt đề tài Thực tập 1 / Thực tập 2}
\gbar{13}{16}{17}{cT3}

% ========================================================
\taskrow{14}{TASK 4 --- Chuẩn bị dữ liệu \& môi trường}
\act{15}{Tải mã nguồn, annotation và checkpoint ProCIR}
\gbar{15}{14}{15}{cT4}
\act{16}{Tìm nguồn ảnh DeepFashion, Fashion200K, FashionGen}
\gbar{16}{15}{17}{cT4}
\act{17}{Trích xuất, sắp xếp ảnh theo product\_id}
\gbar{17}{16}{18}{cT4}
\act{18}{Kiểm chứng độ phủ dữ liệu ảnh}
\gbar{18}{17}{18}{cT4}
\act{19}{Chạy thử trên CPU, đánh giá thời gian chạy}
\gbar{19}{18}{19}{cT4}
\act{20}{Thiết lập môi trường Kaggle GPU}
\gbar{20}{18}{20}{cT4}

% ========================================================
\taskrow{21}{TASK 5 --- Thực nghiệm}
\act{22}{Tái lập ProCIR trên DeepFashion}
\gbar{22}{19}{22}{cT5}
\act{23}{Tái lập ProCIR trên Fashion200K}
\gbar{23}{20}{23}{cT5}
\act{24}{Đối chiếu kết quả tái lập với bài báo}
\gbar{24}{22}{24}{cT5}
\act{25}{Cài đặt baseline CLIP / FashionCLIP (late-fusion)}
\gbar{25}{20}{22}{cT5}
\act{26}{Chạy baseline trên DeepFashion và Fashion200K}
\gbar{26}{22}{24}{cT5}
\act{27}{So sánh ProCIR với các baseline}
\gbar{27}{23}{25}{cT5}

% ========================================================
\taskrow{28}{TASK 6 --- Viết báo cáo}
\act{29}{Dịch toàn văn bài báo sang tiếng Việt}
\gbar{29}{15}{20}{cT6}
\act{30}{Viết chương Giới thiệu và Tổng quan bài báo}
\gbar{30}{18}{22}{cT6}
\act{31}{Viết chương Tái lập thực nghiệm}
\gbar{31}{22}{24}{cT6}
\act{32}{Viết chương Baseline và Kết luận}
\gbar{32}{24}{25}{cT6}
\act{33}{Hoàn thiện định dạng, rà soát và nộp báo cáo}
\gbar{33}{25}{26}{cT6}
\milestone{33}{26}

% ==== Chú thích ====
\fill[cMS] (-\lw+0.5,-\nr-0.5) -- (-\lw+0.9,-\nr-0.9) -- (-\lw+0.5,-\nr-1.3) -- (-\lw+0.1,-\nr-0.9) -- cycle;
\node[anchor=west,font=\sffamily\footnotesize] at (-\lw+1.2,-\nr-0.9) {Mốc nộp báo cáo Thực tập 1 (03/10)};

\end{tikzpicture}
\end{document}
